\chapter{Approach}
\label{chap:approach}

This chapter describes what is computed and how. Section~\ref{sec:preliminaries}
fixes the notation, Section~\ref{sec:architecture} explains the shape of the
pipeline, and Sections~\ref{sec:class_level_metrics}--\ref{sec:global_metrics}
define the thirteen metrics grouped by the level of the graph they describe.
Section~\ref{sec:entropy_weighted_pagerank} presents the entropy-weighted
PageRank variant proposed by this thesis, and
Section~\ref{sec:efficient_computation} covers the techniques that make the
computation tractable on billion-triple graphs.

\section{Preliminaries}
\label{sec:preliminaries}

A knowledge graph is a set of RDF triples $T \subseteq (I \cup B) \times I
\times (I \cup B \cup L)$, where $I$ are \acp{IRI}, $B$ blank nodes and $L$
literals. A triple $(s, p, o)$ states that subject $s$ stands in relation $p$ to
object $o$. Throughout this thesis the following symbols are used:

\begin{itemize}
  \item $T$ --- the set of triples of one snapshot, $|T|$ its cardinality;
  \item $S$ --- the set of distinct subjects, $O$ the set of distinct objects;
  \item $P$ --- the set of distinct predicates;
  \item $C$ --- the set of classes, that is the objects of \texttt{rdf:type} triples;
  \item $E$ --- the set of typed entities, $E_t \subseteq E$ those of class $t$.
\end{itemize}

An entity may carry several types, so the sets $E_t$ overlap and
$\sum_{t \in C} |E_t| \geq |E|$ in general. This matters when reading
Section~\ref{sec:class_level_metrics}: class shares are fractions of $|E|$ but
need not sum to one.

Two versions of the same knowledge graph are written $v_1$ and $v_2$ with triple
sets $T_1$ and $T_2$; two different knowledge graphs are written $A$ and $B$.
The distinction is not cosmetic. Two versions of one graph share their \ac{IRI}
namespace, so their triples can be compared directly; two different graphs
generally do not, and Section~\ref{sec:global_metrics} treats them with separate
metrics for exactly that reason.

All metrics are computed over data served by QLever~\cite{bast2017qlever}, a
SPARQL engine that answers queries from a precomputed index. Every metric in
this chapter is therefore expressible as one or more SPARQL queries plus a
formula applied to their results.

\section{Architecture}
\label{sec:architecture}

The pipeline follows the pattern established by Knowgly~\cite{garcia2025information}:
\emph{SPARQL performs the counting, the host language holds the aggregated counts
in dictionaries and applies the formula, and a separate layer visualises the
result}. Concretely it has four stages:

\begin{enumerate}
  \item a QLever index of one snapshot of a knowledge graph, served as a SPARQL
        endpoint over HTTP;
  \item a set of counting queries, one or a few per metric, evaluated on the
        server so that no triple ever crosses the network unnecessarily;
  \item a Rust program that collects the counts into hash maps, applies the
        metric formula, and writes the result as JSON and CSV;
  \item a niceGUI dashboard that reads those files and renders them.
\end{enumerate}

Splitting the work this way is what keeps the approach applicable to graphs of
this size. The expensive part of every metric is an aggregation over hundreds of
millions of triples, which the engine can answer from its index; what remains is
arithmetic over a few hundred or a few thousand numbers, which is trivial in any
language. The role of Rust is therefore not raw speed in the formula itself but
predictable memory behaviour in the two places where data volume does reach the
client: the triple diff of Section~\ref{subsec:integer_encoding} and the
PageRank iteration of Section~\ref{sec:entropy_weighted_pagerank}.

The endpoint is not compiled in. It is read from the environment, so the same
binary runs against a laptop, a server on the local network, or a university
virtual machine without modification. A second endpoint may be supplied for the
metrics that compare two graphs. Each run may additionally be given a label
naming the snapshot being measured; results are then written to a directory of
that name, which is the mechanism the trajectory metric of
Section~\ref{sec:global_metrics} uses to find several versions of the same
measurement.

\section{Class-Level Metrics}
\label{sec:class_level_metrics}

Class-level metrics describe what kinds of things a graph contains and which
predicates characterise each kind.

\subsection{Class Population and Share}

The most basic question about a snapshot is what is in it. For a class $t$,

\begin{equation}
  \mathrm{share}(t) = \frac{|E_t|}{|E|}
  \label{eq:share}
\end{equation}

gives the fraction of all typed entities belonging to $t$. Both quantities come
from a single grouped count. Tracked across versions, the share exposes which
classes grow, shrink or vanish, and it normalises for overall growth: a class
can double in absolute size while losing share, which is a different fact about
the graph than simply getting bigger.

\subsection{Property Entropy per Type}

Population says nothing about how well a class is described. For that we ask, for
one class $t$ and one predicate $p$, how varied the values of $p$ are. Let $P(o)$
be the relative frequency of value $o$ among all values $p$ takes on entities of
$t$. The Shannon entropy~\cite{shannon1948mathematical} of that distribution is

\begin{equation}
  \mathrm{EntF}(p, t) = - \sum_{o} P(o) \log_2 P(o).
  \label{eq:entf}
\end{equation}

A predicate whose values are spread over many distinct objects --- a birth date,
where nearly every person differs --- scores high and carries real information. A
predicate dominated by one value --- a gender field with two outcomes, or a
constant marker --- scores near zero. Entropy is measured in bits and is bounded
above by $\log_2 D$, where $D$ is the number of distinct values, with equality
when all values are equally frequent.

Computing Equation~\ref{eq:entf} naively means retrieving every object value of
every predicate of every class, which is prohibitive at this scale.
Section~\ref{subsec:count_of_counts} shows how to obtain it exactly from a
result of a few hundred rows.

\subsection{Entropy-Weighted Type Importance}

Property entropy measures how informative a predicate's values are, but not
whether the predicate is characteristic of the class. The complementary signal is
the \ac{ETImp} measure of Knowgly~\cite{garcia2025information}, a TF--IDF
analogue for predicates:

\begin{equation}
  \mathrm{ETImp}(p, t) = \mathrm{EF}_p(p, t) \cdot
    \log_2 \frac{|E_t|}{\mathrm{EF}_p(p, t)},
  \label{eq:etimp}
\end{equation}

where $\mathrm{EF}_p(p,t)$ is the number of distinct entities of class $t$ that
use $p$. A predicate used by every entity of the class contributes
$\log_2 1 = 0$: universal predicates are uninformative precisely because they do
not discriminate.

Neither signal suffices alone. A predicate every entity carries scores zero on
\ac{ETImp} regardless of how diverse its values are; a predicate with a single
repeated value scores near zero on \ac{EntF} regardless of how characteristic it
is. This thesis therefore combines them as a weighted geometric mean:

\begin{equation}
  \mathrm{EntETImp}(p, t) = \mathrm{EntF}(p, t)^{w} \cdot
                            \mathrm{ETImp}(p, t)^{1 - w},
  \qquad w = 0.75.
  \label{eq:entetimp}
\end{equation}

The geometric mean is chosen deliberately over an arithmetic one: it collapses to
zero whenever either factor is zero, which is the desired semantics here. A
predicate must be both typical of the class and informative in its values to rank
highly. The weight $w$ favours the entropy term; it is a parameter, and
Chapter~\ref{chap:experiments} reports its effect.

\subsection{Class Entropy}

The two preceding metrics are defined per predicate. Aggregating to the class
gives one number describing how diverse a class is as a whole:

\begin{equation}
  H(t) = - \sum_{o} P(o) \log_2 P(o),
  \label{eq:class_entropy}
\end{equation}

where $P(o)$ now ranges over \emph{all} object values occurring on entities of
$t$, irrespective of predicate. A class whose entities all look alike scores low;
a class of varied, richly described entities scores high.

The same quantity is also computed over the predicates used by the class rather
than the object values. The two are reported together because they can diverge in
an informative way: a class may take highly varied values through a very narrow
set of predicates, which says something quite different about the modelling than
a class described by many predicates with repetitive values.

\section{Entity-Level Metrics}
\label{sec:entity_level_metrics}

Entity-level metrics move from kinds of things to individual things.

\subsection{Entity Informativeness}

Following the intuition behind InfoRank~\cite{menendez2018ranking} --- that
important entities have a lot recorded about them --- the informativeness of an
entity $v$ is the normalised count of its datatype properties:

\begin{equation}
  \mathrm{IR}(v) = \frac{\mathrm{dtp}(v)}{\sum_{u} \mathrm{dtp}(u)},
  \label{eq:inforank}
\end{equation}

where $\mathrm{dtp}(v)$ is the number of triples of $v$ whose object is a
literal. Literals are what is counted, and the restriction is deliberate: a link
to another entity records who $v$ is related to, whereas a literal records
something \emph{about} $v$. The measure needs a single grouped aggregation and no
iteration, which makes it one of the cheapest metrics in the catalogue. Across
versions it distinguishes entities that are being actively enriched from those
that have been abandoned.

\subsection{Object Diversity}

Entropy weights values by how often they occur. The complementary measure counts
them without weighting:

\begin{equation}
  \mathrm{OD}(p, t) = |\{\, o : (e, p, o) \in T \text{ and } e \in E_t \,\}|,
  \label{eq:objdiv}
\end{equation}

the number of distinct objects a predicate takes within a class. Reported
alongside the total number of such triples, the ratio
$\mathrm{OD}(p,t) / |\{(e,p,o) \in T : e \in E_t\}|$ states whether values are
mostly unique (ratio near one) or heavily reused (ratio near zero). This is an
exact count rather than an information-theoretic quantity, and it serves as a
sanity check on the entropy metrics: the two should agree on the extremes and may
legitimately disagree in between.

\section{Global Metrics}
\label{sec:global_metrics}

Global metrics describe the graph as a whole and, in most cases, the difference
between two graphs.

\subsection{Graph Size and Shape}

The foundation table of every snapshot: the counts $|T|$, $|S|$, $|O|$, $|P|$,
$|C|$ and $|E|$, together with

\begin{equation}
  \mathrm{density} = \frac{|T|}{|E|},
  \qquad
  \overline{\deg}^{\,+} = \frac{|T|}{|S|},
  \qquad
  \overline{\deg}^{\,-} = \frac{|T|}{|O|}.
  \label{eq:shape}
\end{equation}

Density is triples per entity; the two degree averages are the mean out- and
in-degree. None of these is surprising in isolation. Their value is comparative:
computed for every version and every graph, their deltas form the growth table
against which all other comparisons are read. A rise in absolute triple count
means one thing if density rose with it and quite another if density fell.

\subsection{Triple-Level Diff}

The core of any comparison between two versions is the set difference between
their triples:

\begin{equation}
  \mathrm{Added} = T_2 \setminus T_1,
  \qquad
  \mathrm{Deleted} = T_1 \setminus T_2,
  \qquad
  \mathrm{Unchanged} = T_1 \cap T_2.
  \label{eq:diff}
\end{equation}

Conceptually this is elementary. Making it computable over graphs of a billion
triples, through a SPARQL endpoint rather than over local dumps, is not; that is
the subject of Sections~\ref{subsec:integer_encoding}
and~\ref{subsec:subject_window}. This metric applies only to two versions of the
same graph. Two different knowledge graphs generally use disjoint \ac{IRI}
namespaces, and Equation~\ref{eq:diff} evaluated across them reports every triple
of one side as added --- arithmetically correct and entirely without meaning.

\subsection{Class-Level Change Rate}

The diff says how much changed; it does not say where. For each class $t$, let
$A_t$ and $D_t$ be the added and deleted triples whose subject belongs to $t$,
and $T_t(v_1)$ its triples in the older version. The churn of the class is

\begin{equation}
  \mathrm{churn}(t) = \frac{|A_t| + |D_t|}{|T_t(v_1)|}.
  \label{eq:churn}
\end{equation}

This distinguishes two situations a global diff conflates: a graph in which
everything drifted slightly, and a graph in which a few classes were rewritten
while the rest stood still. Because it reuses the diff machinery with the scope
restricted to one class, it costs little beyond the diff itself.

\subsection{Vocabulary and Schema Evolution}

At the schema level the same set difference applies to the class and predicate
vocabularies:

\begin{equation}
  C_{\mathrm{added}} = C_2 \setminus C_1,
  \qquad
  C_{\mathrm{removed}} = C_1 \setminus C_2,
  \label{eq:vocab}
\end{equation}

and analogously for $P$. Presence in the vocabulary is, however, a weak signal on
its own: a term can be introduced and remain almost unused for a version or two.
Each newly added predicate is therefore reported together with the number of
triples that actually use it in $v_2$, which measures adoption rather than mere
declaration.

\subsection{Metric Trajectories}

The framework layer, and the reason every metric writes its result to a labelled
directory. For a metric $m$ and consecutive versions $v_i$,

\begin{equation}
  \Delta m(v_i) = m(v_{i+1}) - m(v_i).
  \label{eq:trajectory}
\end{equation}

The implementation is deliberately generic over $m$: it reads the stored results
of a metric for a list of labelled snapshots, flattens each to its numeric
leaves, and reports the series and its deltas for every value present in all
snapshots. Values that appear in only some snapshots are counted and reported
separately, since their appearance or disappearance is itself a finding. A metric
added later requires no change to this component.

\subsection{Cross-KG Class Comparison}

Finally, the same class may exist in two different knowledge graphs, and its
metric values can be placed side by side:

\begin{equation}
  \Delta m(t) = m_A(t) - m_B(t).
  \label{eq:crosskg}
\end{equation}

Classes are matched at the schema level, by local name. This is a real
limitation and is stated as such: aligning the \emph{instances} of two knowledge
graphs through \texttt{owl:sameAs} links is impractical at this scale, so no
union, overlap or merge gain is estimated here, unlike the analysis of Ringler
and Paulheim~\cite{ringler2017one}. What the comparison does show is how
differently two graphs model and populate the same concept: how many entities
they assign to it, how many predicates they use to describe it, and how diverse
the resulting values are.

\section{Entropy-Weighted PageRank}
\label{sec:entropy_weighted_pagerank}

Classic PageRank~\cite{page1999pagerank} treats every link alike: a node
distributes its rank uniformly over its outgoing edges. In a knowledge graph that
assumption is questionable, because the edges are typed and the types differ
enormously in what they tell us. An edge labelled with a near-constant predicate
connects two nodes without saying much about either; an edge labelled with a
highly varied predicate is far more specific. Knowgly~\cite{garcia2025information}
makes a related observation and weights its links using an InfoRank-derived
factor.

This thesis proposes a simpler variant: weight each edge by the property entropy
of its predicate, the quantity already defined in Equation~\ref{eq:entf}. Let

\begin{equation}
  w(p) = \frac{\mathrm{EntF}(p)}{\max_{q \in P} \mathrm{EntF}(q)}
  \label{eq:weight}
\end{equation}

be the entropy of predicate $p$ over the whole graph, normalised into $[0, 1]$ by
the largest entropy of any predicate. The rank update becomes

\begin{equation}
  \mathrm{PR}(v) = (1 - d) + d \sum_{u \rightarrow v}
                   w(p_{uv}) \cdot \frac{\mathrm{PR}(u)}{\mathrm{outdeg}(u)},
  \label{eq:entropy_pagerank}
\end{equation}

with damping factor $d = 0.85$. Rank now flows preferentially through
informative predicates. Setting $w(p) = 1$ for all $p$ recovers classic PageRank
exactly, which is convenient: the baseline is not a separate implementation but
the same code path with uniform weights, so the two rankings are computed over an
identical edge set and are directly comparable.

Three properties of Equation~\ref{eq:entropy_pagerank} deserve comment. First,
the weight is a property of the predicate and is computed over the entire graph,
not over whatever subgraph the ranking is run on; it is a statement about what
the predicate means in this knowledge graph, not about the sample. Second, only
edges between entities participate --- triples with literal objects are not
links, and including them would let rank leak into values. Third, the
normalisation divides by $\mathrm{outdeg}(u)$, the number of outgoing edges,
rather than by the sum of outgoing weights. This means a node whose outgoing
edges are all low-entropy passes on less total rank than it received, so the
weighted variant is not rank-preserving in the way classic PageRank is. This is a
deliberate consequence: a node that connects to the rest of the graph only
through uninformative predicates should not confer as much importance. It is also
a clear point of comparison against a weight-normalised alternative, and
Chapter~\ref{chap:experiments} reports the difference between the two rankings
rather than either in isolation.

\section{Efficient Computation}
\label{sec:efficient_computation}

The definitions above are straightforward. Evaluating them over graphs of one to
two billion triples, through a query endpoint, on commodity hardware, is where
the work lies. Three techniques carry most of the weight.

\subsection{Entropy from Count-of-Counts Histograms}
\label{subsec:count_of_counts}

A direct evaluation of Equation~\ref{eq:entf} requires the frequency of every
distinct object value of a predicate within a class. For a predicate such as a
label on a class with millions of members, that is millions of rows transferred
and counted --- per predicate, per class.

The observation that removes this cost is that entropy depends only on the
\emph{multiset} of counts, never on the values that carry them. Two predicates
whose value distributions are permutations of one another have identical entropy.
It is therefore sufficient to know, for each multiplicity $c$, how many distinct
values occur exactly $c$ times. Writing $m_c$ for that number, and
$N = \sum_c m_c \, c$ for the total number of value occurrences, the entropy of
Equation~\ref{eq:entf} can be rewritten as

\begin{equation}
  H = \log_2 N - \frac{1}{N} \sum_{c} m_c \cdot c \cdot \log_2 c.
  \label{eq:entropy_histogram}
\end{equation}

\begin{proof}[Derivation]
With $P(o) = c_o / N$ for a value occurring $c_o$ times,
\[
  H = -\sum_{o} \frac{c_o}{N} \log_2 \frac{c_o}{N}
    = \log_2 N - \frac{1}{N} \sum_{o} c_o \log_2 c_o ,
\]
and grouping the values by their multiplicity $c$ replaces the sum over values
by a sum over multiplicities weighted by $m_c$.
\end{proof}

The histogram $(c, m_c)$ is obtained from a nested grouped query: the inner query
counts occurrences per value, the outer counts values per occurrence count. Both
aggregations run on the server, and the result has as many rows as there are
distinct multiplicities --- typically a few hundred, sometimes fewer than ten.
The number of distinct values, needed for Equation~\ref{eq:objdiv}, falls out of
the same histogram as $\sum_c m_c$.

The essential point is that this is not an approximation. Equation~\ref{eq:entropy_histogram}
is an algebraic identity, so the entropy obtained from the histogram is exactly
the entropy of the full distribution. The saving is in transport and client-side
work, not in fidelity, which distinguishes it from sampling-based estimators.

\subsection{Integer-Encoded Triple Sets}
\label{subsec:integer_encoding}

The set operations of Equation~\ref{eq:diff} are the one place where triples
themselves, rather than counts, must reach the client. Comparing them as strings
is expensive in both time and space: each triple carries three \acp{IRI} of
typically 30 to 60 bytes, and every set operation hashes all three.

Instead, each distinct term encountered on either side is interned once into a
32-bit integer, with a single dictionary shared across both versions so that
equal strings receive equal identifiers. A triple becomes a triple of integers
occupying twelve bytes, and the comparison reduces to sorting two integer arrays
and merging them linearly. String data is touched exactly once, during
interning; the set operations themselves never look at a byte of text.

Both this path and a naive implementation using hash sets of string triples are
implemented, and Chapter~\ref{chap:experiments} benchmarks them against each
other. The two are also cross-checked at run time: the implementation asserts
that both produce identical added, deleted and unchanged counts, so the reported
speedup cannot silently be a speedup of the wrong answer.

\subsection{Subject-Window Scoping}
\label{subsec:subject_window}

Even encoded, the full triple set of a billion-triple version does not fit in the
memory of the machines used here, so the diff must run over a well-defined
subset. The subset has to satisfy two conditions: it must be cheap to obtain, and
both versions must return \emph{the same} subset, or the difference between them
is an artefact of the sampling rather than of the data.

A bare \texttt{LIMIT} fails the second condition, since the engine may return any
rows it likes. The obvious repair, ordering the triples and taking a prefix,
fails the first: ordering by subject, predicate and object forces the server to
sort the entire graph before it can return the first row. Measured on a graph of
roughly 1.5 billion triples, QLever attempted to allocate 12.7\,GB for such a
query and aborted.

The scheme adopted instead is a \emph{subject window}. Ordering a list of
distinct subjects is far cheaper than ordering every triple, so from each version
the lexicographically first $n$ subjects are retrieved. The two prefixes are then
truncated at whichever ends earlier, and their union forms the window. Every
triple of every subject in the window is fetched by naming the subjects
explicitly, which requires no ordering at all, since naming them already
determines the result set exactly.

The window has a property worth stating precisely: within it, both versions are
covered \emph{completely}. Any triple of any subject in the window is present in
the fetched data for whichever versions contain it, so the diff restricted to the
window is exact --- including for entities that appear in one version and not the
other. Truncating at the earlier of the two upper bounds is what guarantees this;
beyond that bound one version would be missing subjects the other has, and their
absence would be indistinguishable from deletion. The diff is thus not an
estimate of the global diff, but an exact diff of a precisely characterised
region, and the region can be enlarged by increasing $n$ or restricted to a class
of interest.

\subsection{Query Formulation}
\label{subsec:query_formulation}

A final class of savings comes from formulating the counting queries so that the
engine does no work the semantics do not require. The clearest example concerns
counting the members of each class. The natural formulation counts distinct
subjects per type, but under the set semantics of RDF the subject--type pairs are
already unique, so a plain count returns exactly the same numbers while sparing
the engine a deduplicating sort. On the graphs used here the equivalent query is
approximately 1.6 times faster, and the equality of the two results is verified
rather than assumed.

The complementary case is worth stating too, since the substitution is not
universally valid: counting distinct entities that \emph{use} a given predicate
does require deduplication, because an entity may use the same predicate several
times. Equation~\ref{eq:etimp} depends on that distinction, and the
implementation preserves it.

\section{The Dashboard}
\label{sec:dashboard}

The visual-analytics component reads the stored metric results and renders them.
No SPARQL query is issued when a page is viewed: the metrics are precomputed and
the dashboard reads their JSON output, so interaction is immediate and
independent of the availability or load of the endpoint. The cost of this
decision is that the dashboard shows the state of the last computation rather
than of the live graph, which is the right trade for metrics that describe
published snapshots and change only when a new version is released.

Because results are stored per labelled snapshot, the same view can be populated
from several versions at once, which is what turns the per-snapshot metrics of
this chapter into the evolution views that
Chapter~\ref{chap:experiments} evaluates.
